\chapter{A product identity repairs the endpoint at every coefficient rate}
\label{component:p05180564ed05}
\noindent\textbf{Component }\path{p05180564ed05}\par
{\small\noindent Source: \path{CORE/prior/inherited/product_repair.tex}\par}
{\small\noindent SHA256: \path{05180564ed0508f8b2ad17b0027feb9f0cbd47ec865e82f3ee046c30c5b545cf}\par}
{\small\noindent Inventory use: actual bounded-b W-state nine-root proof\par}
\medskip
\section{A product identity repairs the endpoint at every coefficient rate}\label{p05180564ed05:sec:productrepair}
The new boundary mechanism is algebraic. A troublesome root is retained, and a second relation controls its product with an additional root. This supplies two positive pivots without comparing an infinitesimal in one field with an infinitesimal in the other.

\begin{lemma}[Product-relation completion]\label{p05180564ed05:lem:productrepair}
Let $z$ be a root vector over $T$, and $h,r\in T$. Suppose
\[
F=(z,h,r)^TQ(z,h,r),\qquad Q\in\Sym(B),\qquad
\bar Q=\operatorname{diag}(Q_0,-a,0),
\]
where $Q_0$ is a real positive-definite matrix and $a\in\R_{\ge0}$. Suppose $p\in B$ has nonzero residue, $A$ is a row with entries in $\mathfrak n$, and
\[
ph=Az+r,\qquad phr=(z,h,r)^TE(z,h,r),\quad E\in\Sym(\mathfrak n).
\]
Then $F\in\Sigma T^2$.
\end{lemma}
\begin{proof}
For a sufficiently large positive real $M$, add the exact zero
\[
M\big[p^2h^2+r^2-(Az)^2-2phr\big]=0.
\]
The residue Gram becomes $\operatorname{diag}(Q_0,M\bar p^2-a,M)\succ0$. The subtracted terms have zero residue by hypothesis. Lemma~\ref{p25af2c58da7e:lem:gram} applies. In particular $p$ is never inverted.
\end{proof}

\subsection{Exact contact scope and full-polynomial input}
This application uses the canonical C1 contact model, with $\lambda=n=0$. Its source parameters are positive $e,k,b,a,t$ satisfying
\[
e,k\in\mathfrak m_c,\quad b\in B_c,\qquad
R_c=1+k^2(1-ae^2b^2),\quad T_0=1+k^2-tk^4b^2,
\]
\[
0<a<T_0,\qquad0<t<R_c.
\]
The signed weight satisfies $d>0$, and all eight event coefficients are nonnegative original physical C1 contact coefficients. This is a specified contact chart, not every coefficient vector with the same omissions. The independent testing chart is $W=1$ with $u,v,X,Y,Z$ infinitesimal. In this section there is no restriction on $k^3/e$ or $k/e$.

The complete seven-packet contact polynomial is denoted $G$. Its short factored definition and unchanged original coefficient vector are printed below. The accompanying checkers compare the entire polynomials, including the scalar and normal terms; no global coverage claim is inferred from those checks.

Define
\begin{gather*}
v_g=ekb^2u,\quad X_g=-uZ+tv_g,\quad Y_g=(1-ae^2b^2)u,\\
N_1=X-X_g,\quad N_2=Y_g-Y,\quad \delta v=v-v_g,\quad D_t=Z-uX,\\
U_0=1+u^2,\quad c_0=1-tk^2b^2,\quad P_0=T_0+u^2=U_0+k^2c_0,\\
S=u+uvX-Y-vZ+(T_0-a)e^2b^2u,\\
s_0=eb^2u(e-kZ),\qquad q=eb^2(eu^2-k^3Z).
\end{gather*}
\par\noindent\textit{Source inclusion. }The original input is reproduced once as Component p7d4e5164e84a (Chapter~\ref{component:p7d4e5164e84a}).\par
Direct substitution gives the exact nonlinear relation
\[
L=P_0s_0-S+N_2-D_t\delta v+uv_gN_1-uc_0q=0. \tag{E1}
\]
Peel the two complete current squares:
\[
E=G-H_3(1+v^2)N_2^2-H_5S^2,\qquad
H_3=4R_cc_3,\quad H_5=4R_cc_5,\quad \bar H_3=\bar H_5=8.
\]
Each peeled packet divided by $e^2$ is polynomial in bounded variables. The full graph scalar has the identity
\[
E_g=A_{00}s_0^2+A_{01}s_0q+A_{11}q^2,
\]
and the entire first normals, with $g_0=s_0/(eb^2)$ and $g_2=q/(eb^2)$ denoting their displayed polynomial expressions, satisfy
\[
E_X=e^2b^2(P_Xg_0+Q_Xg_2),\quad
E_Y=eb^2(P_Yg_0+Q_Yg_2),\quad
E_v=e^2b^2(P_vg_0+Q_vg_2).
\]
The exact polynomial splits and residues needed are
\begin{center}\small
\begin{tabular}{ccc}\toprule
Quantity&Polynomial split&Quotient residues\\\midrule
$A_{00}$&$k^4a_0+e^2b^2a_1$&$-8,-8$\\
$A_{01}$&$k^2c_1+ebc_2$&$0,0$\\
$P_X$&$k^2p_X+ebq_X$&$0,0$\\
$P_Y$&$k^2p_Y+ebq_Y$&$0,0$\\
$P_v-16kP_0$&$k^2p_v+ebq_v$&$0,0$\\\bottomrule
\end{tabular}
\end{center}
Also $\bar A_{11}=8$ and $\bar Q_X=\bar Q_Y=\bar Q_v=0$. These equalities remain valid with $b$ symbolic; boundedness, rather than infinitesimality, suffices for $b$.

\par\noindent\textit{Source inclusion. }The original input is reproduced once as Component p935c9d6318ea (Chapter~\ref{component:p935c9d6318ea}).\par

Exact finite Taylor expansion in $(N_1,-N_2,\delta v)$ accounts for every higher-normal term. Because $E/e^2$ is polynomial, every term of normal degree at least two allocates to $x=eN_1$, $y=e\delta v$, $N_2$ with bounded coefficients. Terms of degree at least three have an extra infinitesimal normal. Terms involving $N_2$ retain an $e$ factor. The only additional nonzero residue is
\[
B_0=\begin{pmatrix}16&-8\\-8&8\end{pmatrix}
\quad\text{on }(x,y).
\]
It is visible without expansion of the full Hessian: at $u=Y=Z=0$, the quadratic residue of $E/e^2$ is $8(X-v)^2+8X^2$, with $b$ still free.

\subsection{The new identity and its nine-root proof}
If the signed testing core is nonnegative, its pure testing square root and the original event packets already give SOS. Otherwise the sign gauges $u,v\ge0$ preserve every packet, and the core Schur complement implies $u,v>0$ and $D_t>0$. To check the latter, after $Z=D_t+uX$ the Schur complement on $(W,D_t)$ is
\[
\frac{2(1+v^2)}{N_0}
\begin{pmatrix}
u^2v^2(1+u^2)&-uv(1+u^2)(1+v^2)\\
-uv(1+u^2)(1+v^2)&u^2v^2
\end{pmatrix},
\]
where $N_0=u^2v^4+u^2v^2+u^2+v^4+2v^2+1>0$. The eliminated block is positive definite. For $W=1,D_t\le0$, the remaining form is nonnegative, proving the claim.

Consequently $w=\sqrt{ek}\sqrt{D_t}\,\delta v$ is a legal separated root. Set
\begin{gather*}
h=k^2s_0-uq=e^2b^2u(k^2-u^2),\qquad z=ebs_0,\qquad r=k^2D_t\delta v,\\
A_0=k^2(S-N_2)-k^3b^2u^2x-uU_0q.
\end{gather*}
Multiplication of (E1) by $k^2$ gives $P_0h=A_0+r$. The additional exact identity is
\begin{align*}
P_0hr={}&ku^2(P_0-k^2)y(S-N_2)
-ku(P_0+c_0k^2u^2)yq\\
&+u^2(P_0-k^2)w^2-k^4b^2u^2T_0xy. \tag{E2}
\end{align*}
Every displayed coefficient has zero residue and is a polynomial in separately bounded factors. To derive (E2), substitute the two soft roots to obtain
\[
u(P_0-k^2)s_0-(1+c_0u^2)q
=ekb^2(k^2-u^2)(D_t+uN_1).
\]
Then use $D_t=U_0Z-uN_1-tu v_g$, $T_0=1+k^2c_0$, and (E1) to replace $S-N_2$ in (E2). The terms in $D_t(\delta v)^2$ cancel; the remaining terms give $P_0hr$. This is a whole-polynomial calculation.

Let $J_0$ denote the right-hand side of (E2). Thus
\[
Z_0=P_0^2h^2+r^2-A_0^2-2J_0=0.
\]
Add to the entire $G$ the exact zero expression
\[
-16kyL+16(ebL)^2+16Z_0.
\]
The first term cancels the exceptional $16kP_0$ first normal and supplies $16w^2$. Use $k^2s_0=h+uq$ in every other scalar and first-normal term. The scalar becomes
\[
a_0(h+uq)^2+a_1z^2+[c_1(h+uq)+c_2z]q+A_{11}q^2.
\]
Its only negative residue diagonals are $-8$ on $h,z$. The safe relation
\[
ebL=P_0z-ebS+ebN_2-bD_ty+ekb^3u^2x-ebu c_0q
\]
has only infinitesimal coefficients apart from $P_0$; hence $16(ebL)^2$ changes the $z$ diagonal to $8$. Every coefficient of $A_0$ is infinitesimal, as is every coefficient in (E2). Thus $16Z_0$ changes the $h$ diagonal to $8$ and supplies a new $r$ diagonal $16$, with no other residue change.

The full Gram on
\[
\boxed{(x,y,N_2,S,w,q,h,z,r)}
\]
therefore has residue
\[
\boxed{B_0\oplus\operatorname{diag}(8,8,16,8,8,8,16)\succ0.}
\]
Its first two LDL pivots are $16,4$; the other seven pivots are printed above. The strict-residue lemma proves SOS of the complete seven-packet polynomial for every coefficient ratio $k^3/e$.

\subsection{Restoration and coverage}
The complete original-parent restoration is
\begin{align*}
\Lambda F={}&G
+4T_0e^2c_2\frac{t}{M}(r_2-Mq_0)^2
+4T_0e^2c_\beta\frac{M}{t}(1+v^2)(r_\beta-tq_0)^2\\
&+4R_cc_3\frac{J}{a}(1+v^2)(r_3-az_0)^2
+4R_cc_5\frac{a}{J}(r_5+Jz_0)^2.
\end{align*}
Every additional packet is a nonnegative coefficient factor times a pure-testing factor times a square. This identity is four instances of
\[
S^2+2(A+B)qS+(A+B)Aq^2+\frac BA S^2
=\frac{A+B}{A}(S+Aq)^2.
\]
Use respectively $(A,B,q)=(M,t,q_0),(t,M,q_0),(a,J,z_0),(J,a,-z_0)$ and $S=r_2-Mq_0,r_\beta-tq_0,r_3-az_0,r_5+Jz_0$. Comparing the seven resulting coefficients with the displayed source proves the identity. Since $\Lambda$ is a positive pure coefficient factor, the pure-factor restoration in \hyperref[component:p9b00698ff36a]{RFA13} (p.~\pageref{component:p9b00698ff36a}) recovers the original parent. Normalizing the original state to $W=1$ and undoing that normalization use only testing factors.

This removes the former $k^3/e$ restriction and includes positive bounded $b$. It does not prove entry of every physical C1 or C2 source into this domain. For unbounded $b$, the recovered inverse-$b$ theorem remains a different input with a ten-root, large polynomial certificate. The new short identity should not be presented as a replacement for that certificate or for the full source cover.
